\chapter{1907: Eduard Buchner}
\label{chap:chemistry1907}

The Nobel Prize in Chemistry for the year 1907 was awarded to Eduard Buchner.

\textbf{Motivation:}
for his biochemical researches and his discovery of cell-free fermentation

\section{Eduard Buchner}

\subsection*{Early Life and Education}

Eduard Buchner was born on 1860-05-20 in Munich, Bavaria (now Germany).
Eduard Buchner was a German
chemist
and expert on fermentation, awarded the 1907
Nobel Prize in Chemistry for his work on
fermentation.

Buchner came from a family with a strong tradition in science: his older brother Hans was a bacteriologist who later became famous for his discovery of complement, a component of the immune system. Eduard himself studied at the University of Munich, where he worked under Adolf von Baeyer, and it was during this period that he formed a lasting friendship with the chemist Theodor Curtius. His doctoral research concerned the synthesis of some nitrogen-containing organic compounds, but his interests were already turning toward the chemistry of living processes, the field in which he would make his name.

\subsection*{Scientific Context}

The nature of fermentation was one of the oldest and most contested questions in chemistry and biology. Since the work of Louis Pasteur in the 1850s and 1860s, it was widely accepted that fermentation, the conversion of sugar into alcohol and carbon dioxide, was brought about by living yeast cells and could not proceed without them. Pasteur himself had declared that fermentation is the result of life without air, and this view, championed by the vitalists, held that the chemical activity of the yeast cell was inseparable from the cell itself. A contrary view had been advanced in 1897 by the Russian chemist Marie Manasseina and argued by others, holding that the active principle could be extracted from the cell and could act on its own, but no one had yet produced convincing evidence for it.

\subsection*{Career and Major Research}

Buchner held professorships in Tübingen, Würzburg, and Berlin. In 1897 he showed that a cell-free extract of yeast, which he named zymase, could ferment sugar to alcohol and carbon dioxide, proving that fermentation is brought about by enzymes rather than by the living yeast cell. This discovery founded the study of enzymes outside the organism.

The discovery arose from an attempt to solve a completely different problem. Buchner's brother Hans had been interested in the preparation of extracts of yeast cells for use in the treatment of disease, and he asked Eduard to help him prepare such extracts by grinding the yeast with sand and pressing the resulting pulp. Eduard Buchner's aim was to obtain a cell-free liquid that might be useful in biochemical work, and to preserve it he added sugar, a common method of keeping such preparations from spoiling. To his astonishment, the sugar was fermented: the cell-free extract produced alcohol and carbon dioxide just as living yeast did. The crucial control, the absence of intact cells, was checked by microscopic examination, and the conclusion was inescapable.

\subsection*{Nobel Prize-winning Work}

The work for which Eduard Buchner was awarded the Nobel Prize in Chemistry in 1907 was described by the Nobel Committee as: "for his biochemical researches and his discovery of cell-free fermentation".

This research fundamentally advanced biochemistry by demonstrating that fermentation could occur in a cell-free extract and by establishing enzymes as chemical catalysts.

The discovery was published in 1897 under the title "Alkoholische Gärung ohne Hefezellen", and it created a sensation. If a cell-free extract could ferment sugar, then the activity of the yeast cell was due not to some vital force but to a chemical substance, which Buchner named zymase. The idea that fermentation was a chemical process, argued by some earlier chemists but abandoned under Pasteur's influence, was suddenly restored with full experimental support, and the study of the enzymes of fermentation was opened to chemists as a branch of ordinary chemistry.

\subsection*{The Work in Detail}

Buchner's experimental method was decisive because it settled the question that had defeated his predecessors: the complete absence of living cells. The yeast was ground with quartz sand and kieselguhr in a hydraulic press, and the clear liquid that emerged was filtered with care and examined microscopically. No cells were found, and the extract nevertheless fermented sugar solutions vigorously, producing alcohol and carbon dioxide in the proportions characteristic of ordinary fermentation. Buchner also showed that the fermenting power of the extract was destroyed by heat and by the addition of antiseptics such as toluene, behavior that would later be recognized as typical of enzymes.

The name zymase was applied to the active principle in the extract, but Buchner was careful to point out that the fermentation of sugar to alcohol is a complex process that may require more than one enzyme. This proved to be correct: the conversion of glucose to ethanol is now known to proceed through many enzymatic steps. Buchner's demonstration that the whole sequence could occur in a cell-free extract removed fermentation from the realm of vital phenomena and placed it firmly within the domain of chemistry, where it could be analyzed step by step.

\subsection*{Reception}

The reaction to Buchner's discovery was immediate and profound. Many biologists were skeptical at first, unwilling to accept that the activities of a living cell could be reproduced by a lifeless extract, but the evidence was repeated and confirmed in laboratories throughout Europe within a few years. Within a decade the new science of enzymology, the study of enzymes as chemical substances, was firmly established. Buchner's discovery is often described as the event that removed the boundary between the chemistry of living things and the chemistry of the laboratory.

\subsection*{Significance and Impact}

The impact of Eduard Buchner's work extends far beyond the specific achievement recognized by the Nobel Prize.
These contributions continue to influence chemical research and education today.

The demonstration that fermentation is a chemical process catalyzed by enzymes opened the way to the modern understanding of metabolism. Once it was accepted that the reactions of living cells are ordinary chemical reactions, they could be studied in the test tube, purified, and ultimately understood one by one. The detailed map of the metabolic pathways, by which sugar is broken down to yield energy in living organisms, was made possible in part by Buchner's demonstration that the process could occur outside the cell, and modern biochemistry rests on this foundation.

\subsection*{Later Career and Honors}

Eduard Buchner continued his scientific work until 1917-08-13, passing away in Focsani, Romania.
He received numerous honors including membership in major scientific academies and societies.

Buchner's later career was spent largely at the universities of Würzburg and Berlin, where he continued his studies of fermentation and of the enzymes involved in it. He received many honors, including honorary doctorates and membership in learned societies, and he was at the height of his reputation when the First World War intervened. He volunteered for military service and was wounded, and in 1917 he died at the field hospital in Focsani, Romania, as a result of complications following a wound received at the front.

\subsection*{Legacy}

The legacy of Eduard Buchner endures through the lasting impact of his scientific contributions.
Eduard Buchner's work continues to be cited and built upon by researchers across the chemical sciences.

Buchner's discovery is commemorated as the birth of enzyme chemistry, and the demonstration that cellular processes could be studied outside the cell is recognized as one of the decisive advances in the history of biology and chemistry. His name is associated with the zymase concept and with the cell-free study of metabolism, and the methods he introduced for grinding cells and preparing active extracts remain in use in modern laboratories, in the form of the cell homogenates and cell-free extracts that are routine tools of biochemical research.

\subsection*{Modern Developments}

Buchner's demonstration of cell-free fermentation proved that at least some biological processes are chemical reactions catalyzed by enzymes, helping to establish biochemistry and enzymology. Today, industrial biocatalysis uses purified enzymes and engineered whole cells for everything from biofuels and food production to the synthesis of pharmaceuticals.

The understanding of fermentation that Buchner made possible has profound practical consequences. The industrial production of ethanol for fuel, beverages, and industrial solvents depends on fermentation carried out on an enormous scale, and the brewing and baking industries are themselves applications of the enzyme chemistry he opened up. Modern biotechnology, in which microorganisms are engineered to produce drugs, enzymes, and materials, and in which isolated enzymes are used to carry out specific chemical transformations in industrial reactors, is the direct descendant of the discovery that a cell-free extract could carry out the chemistry of life.

\subsection*{Selected Publications}

"Alkoholische Gärung ohne Hefezellen" (1897)
"Die Zymasegärung" (1903)
